\section{관련 연구}
\label{sec:related}

\subsection{배틀로얄과 e스포츠의 결과 예측}
\label{sec:related:esports}

% P1: 경기 결과 예측에서 경기 중 위험 예측으로
배틀로얄 게임의 예측 연구는 주로 경기가 끝난 뒤의 통계로 최종 순위를 예측하는 문제에서 출발하였다.
Kaggle의 PUBG 순위 예측 대회는 경기 종료 후의 플레이어 통계로 최종 순위를 회귀하는 과제를 제시하였고~\cite{KagglePUBG}, Lee와 Lee는 같은 유형의 통계로 순위를 예측할 때 유효한 요인과 모델이 경기 모드에 따라 달라짐을 보였다~\cite{LeeLee25}.
배틀로얄의 결과에는 무작위성이 크게 작용한다는 분석도 있다~\cite{Galka21,Rosenbusch20}. 이는 결과를 결정적으로 맞히기보다 확률로 다루어야 함을 시사한다.
다른 장르의 e스포츠에서는 경기 중 예측이 활발히 연구되었다.
MOBA 장르에서는 실시간 승리 예측~\cite{Hodge21}과 시간 구간별 팀 간 상대적 우위의 평가~\cite{Yu19}가 제안되었다.
개인 단위의 짧은 구간 사망 예측도 Dota~2~\cite{Katona19,Ringer23}와 Counter-Strike~\cite{Marshall22}에서 다루어졌다.
배틀로얄의 경기 중 예측으로는 Sagong 등이 단계별 활동 점수와 플레이 유형으로 개별 플레이어의 생존을 예측하였고~\cite{Sagong25a}, 단계별 전략 시퀀스로 승리를 예측하였다~\cite{Sagong25b}.
KRAFTON은 생존분석을 이용하여 팀 단위의 실시간 우승확률을 추정하는 예측기를 공개하였다~\cite{KRAFTON}.
이들 연구는 경기 단위의 결과(순위·승리) 또는 개인 단위의 짧은 구간 사망을 대상으로 한다.
줄어드는 안전구역의 단계를 위험 구간으로 하는 팀 단위의 탈락 위험은 검색한 범위에서 다루어지지 않았으며, 배틀로얄의 기존 설계들을 같은 자료에서 비교한 연구도 찾지 못하였다.

\subsection{생존분석과 이산시간 hazard}
\label{sec:related:survival}

% P2: 추정 대상과 기저위험
생존분석은 사건이 일어나기까지의 시간과 관측 종료를 함께 다룬다.
Cox 비례위험모형은 부분우도로 공변량의 상대위험을 추정하며, 기저위험을 지정하지 않는다~\cite{Cox72}.
반면 이산시간 생존모형은 시간을 구간으로 나누고, 이전 구간까지 생존한 대상이 각 구간 안에 사건을 겪을 조건부 확률, 즉 구간별 hazard를 구간별 기저위험과 함께 추정한다~\cite{SingerWillett93}.
이 구조는 신경망과 결합되어 대규모 자료로 확장되었고~\cite{Gensheimer19}, 생존 시간의 분포를 직접 학습하는 모형~\cite{Lee18}, 시간에 따라 갱신되는 공변량을 이용하는 동적 예측~\cite{Lee20}, 순환 신경망으로 시점별 조건부 사건 확률을 모형화하는 방법~\cite{Ren19}이 제안되었다.
예측 시점마다 그 시점의 생존자를 대상으로 정해진 구간의 위험을 추정하는 landmarking도 동적 예측의 표준적인 방법이다~\cite{vanHouwelingen07}.
생태학의 점유 모형 역시 기간별 국지 멸종 확률을 추정한다는 점에서 같은 이산시간 구조를 가진다~\cite{MacKenzie03}.
본 연구의 다음 단계 탈락 위험은 각 단계 경계를 예측 시점으로, 다음 단계를 예측 구간으로 하는 이산시간 hazard이다.
따라서 Cox 계열 학습 목적으로 얻는 상대위험 점수와는 추정 대상이 다르며, 본 연구는 두 학습 목적을 같은 자료와 같은 모델 구조에서 비교한다.

\subsection{시공간 다중 에이전트 모델링}
\label{sec:related:multiagent}

% P3: 상호작용 구조와 표현
시공간 데이터 마이닝은 이동 궤적과 사건을 이용한 예측, 군집, 이상 탐지 등을 다루어 왔으며, 최근에는 딥러닝이 주요 방법이 되었다~\cite{Atluri18,Wang22}.
다수의 에이전트가 상호작용하는 장면에서는 에이전트 축과 시간 축의 attention을 결합하여 궤적과 상태를 함께 모형화하는 방법이 제안되었고~\cite{Ngiam22,Capellera24}, 그래프 구조를 attention의 편향으로 주입하는 방법도 널리 쓰인다~\cite{Ying21}.
스포츠와 e스포츠에서도 선수 간 상호작용 그래프를 이용하여 경기 결과를 예측한 연구가 있다~\cite{Xenopoulos21}.
이러한 구조는 상호작용을 학습할 수 있는 표현력을 제공하지만, 입력 표현과 학습 목적이 함께 바뀌는 비교에서는 구조 자체의 기여를 분리하기 어렵다.
본 연구는 팀 축 attention, 단계 축 attention, 관계 편향을 갖춘 모델을 참조 모델로 사용하고, 입력 표현과 학습 목적을 고정한 상태에서 각 구조 요소를 하나씩 추가하여 그 기여를 평가한다.
